\subsection{\pp vs. \mpp by Example}
\label{sec:pp_example}

Since the last sections were mainly of a theoretical nature, in this section, I want to compare the \pp and \mpp metrics and highlight their differences and shortcomings using an example. 

\begin{table}[]
    \centering
    \begin{tblr}{colspec = {Q[c,m] Q[c,m] *{4}{Q[r,m]}},
                 row{1} = {font=\bfseries, bg=gray!50},
                 row{odd[2]} = {bg=gray!25},
                 rowsep=4pt,
                 row{1}={c},
                 }
        \toprule
        Platform & {arch.\\class} & {theo. peak\\performance} & App. 1                    & App. 2                   & App. 3                   \\\midrule
        A        & 1              & $\SI{9.75}{\tera\flops}$  & $\SI{7.1}{\tera\flops}$   & $\SI{4.6}{\tera\flops}$  & $\SI{0.4}{\tera\flops}$  \\
        B        & 1              & $\SI{5.16}{\tera\flops}$  & $\SI{3.9}{\tera\flops}$   & $\SI{2.5}{\tera\flops}$  & $\SI{0.16}{\tera\flops}$ \\
        C        & 1              & $\SI{22.63}{\tera\flops}$ & $\SI{16.4}{\tera\flops}$  & $\SI{10.8}{\tera\flops}$ & $\SI{5.3}{\tera\flops}$  \\
        D        & 2              & $\SI{1.65}{\tera\flops}$  & $\SI{0.02}{\tera\flops}$  & $\SI{0.9}{\tera\flops}$  & $\SI{0.02}{\tera\flops}$ \\
        E        & 2              & $\SI{1.32}{\tera\flops}$  & $\SI{0.01}{\tera\flops}$  & $\SI{0.7}{\tera\flops}$  & $\SI{0.03}{\tera\flops}$ \\
        F        & 2              & $\SI{0.61}{\tera\flops}$  & $\SI{0.005}{\tera\flops}$ & ---                      & $\SI{0.01}{\tera\flops}$ \\
        \bottomrule
    \end{tblr}
    \caption{%
    Architectural performance results for three example applications and six hardware platforms. 
    Platforms A, B, and C belong to one architectural class, and platforms D, E, and F belong to another class. 
    Application 1 had a very good performance in architectural class 1 but a very bad performance in class 2. 
    Application 2 performs well across all platforms except for platform F, which is unsupported. 
    Application 3 performs badly on all platforms. 
    }
    \label{tab:example_perfs}
\end{table}

\autoref{tab:example_perfs} shows six different platforms from two different architectural classes, for example, \glspl{gpu} and \glspl{cpu}. 
Column 3 of \autoref{tab:example_perfs} shows the theoretical architectural peak performances of the six platforms in $\si{\tera\flops}$. 
The remaining three columns show three different applications with different characteristics. 
Application one performs great on architectural class one but poorly on class two; application two performs well on all platforms regardless of architectural class, except for platform F, which is completely unsupported; and application three performs very poorly across all platforms and architectural classes. 
These example applications reflect the performance characteristics of real-world applications: application one may be specifically optimized for \glspl{gpu} but not for \glspl{cpu}, application two is optimized for \glspl{gpu} and \glspl{cpu}, but, e.g, some new ARM \gls{cpu} is unsupported, and application three is not optimized whatsoever. 

\begin{table}[]
    \centering
    \begin{tblr}{colspec = {Q[c,m] *{3}{Q[r,m]}},
                 row{1} = {font=\bfseries, bg=gray!50},
                 row{odd[2]} = {bg=gray!25},
                 rowsep=4pt
                 }
        \toprule
        Platform & App. 1        & App. 2        & App. 3 \\\midrule
        A        & $\num{0.73}$  & $\num{0.47}$  & $\num{0.04}$  \\
        B        & $\num{0.76}$  & $\num{0.48}$  & $\num{0.03}$  \\
        C        & $\num{0.72}$  & $\num{0.48}$  & $\num{0.23}$  \\
        D        & $\num{0.01}$  & $\num{0.55}$  & $\num{0.01}$  \\
        E        & $\num{0.008}$ & $\num{0.53}$  & $\num{0.02}$  \\
        F        & $\num{0.008}$ & $\num{0.0}$   & $\num{0.02}$  \\
        \bottomrule
    \end{tblr}
    \caption{%
    The architectural efficiencies for the example applications, hardware platforms, and performance results in \autoref{tab:example_perfs}. 
    The efficiency of application two on platform F is $\num{0.0}$ since the platform is unsupported. 
    }
    \label{tab:example_effs}
\end{table}

Based on the raw performance results on \autoref{tab:example_perfs}, we can calculate the architectural efficiency based on \autoref{def:arch_eff}.
The resulting architectural efficiencies are shown in \autoref{tab:example_effs}. 
As we can see, \autoref{tab:example_effs} reflects the previously discussed application characteristics. 
Note that application two's architectural efficiency on platform F is zero since this platform is unsupported for application two. 

\begin{table}[]
    \centering
    \begin{tblr}{colspec = {Q[l,m] *{7}{Q[c,m]}},
                 row{1,2} = {font=\bfseries, bg=gray!50},
                 row{odd[2]} = {bg=gray!25},
                 rowsep=4pt,
                 cell{1}{1}={r=2}{l},
                 cell{1}{2}={r=2}{c},
                 }
        \toprule
        Platform Sets $\mathbf{H}$     & $HS$      & \SetCell[c=2]{c} App. 1        && \SetCell[c=2]{c} App. 2        && \SetCell[c=2]{c} App. 3        &\\
                                       & $\num{0}$ & \pp           & \mpp            & \pp           & \mpp            & \pp           & \mpp            \\\midrule
        \{ A \}                        & $\num{0}$ & $\num{0.730}$ & $\num{0.730}$   & $\num{0.470}$ & $\num{0.470}$   & $\num{0.040}$ & $\num{0.040}$   \\
        \{ A, B \}                     & $\num{0}$ & $\num{0.745}$ & $\num{0.745}$   & $\num{0.475}$ & $\num{0.475}$   & $\num{0.033}$ & $\num{0.035}$   \\
        \{ A, B, C \}                  & $\num{0}$ & $\num{0.736}$ & $\num{0.737}$   & $\num{0.477}$ & $\num{0.477}$   & $\num{0.048}$ & $\num{0.100}$   \\
        \{ A, B, C, D \}               & $\num{0.811}$ & $\num{0.038}$ & $\num{0.555}$   & $\num{0.493}$ & $\num{0.495}$   & $\num{0.025}$ & $\num{0.078}$   \\
        \{ A, B, C, D, E \}            & $\num{0.971}$ & $\num{0.022}$ & $\num{0.446}$   & $\num{0.500}$ & $\num{0.502}$   & $\num{0.024}$ & $\num{0.066}$   \\
        \{ A, B, C, D, E, F \}         & $\num{1}$ & $\num{0.017}$ & $\num{0.373}$   & $\num{0.0}$   & $\num{0.502}$   & $\num{0.023}$ & $\num{0.058}$   \\\midrule
        \{ A, B, C \} (arch. 1)  & $\num{0}$ & $\num{0.736}$ & $\num{0.737}$   & $\num{0.477}$ & $\num{0.477}$   & $\num{0.048}$ & $\num{0.100}$   \\
        \{ D, E, F \} (arch. 2)  & $\num{0}$ & $\num{0.009}$ & $\num{0.009}$   & $\num{0.0}$   & $\num{0.540}$   & $\num{0.015}$ & $\num{0.017}$   \\
        \bottomrule
    \end{tblr}
    \caption{%
    The performance portability scores of all three applications based on the architectural efficiency values from \autoref{tab:example_effs} for \pp and \mpp on different platform sets adding more and more platforms to the platform set $\mathbf{H}$ also showing the respective proposed heterogeneity score $HS$. 
    Additionally, as \citet{marowka_toward_2021} proposed, separate performance portability scores for the different architectural classes are reported.  
    }
    \label{tab:example_pp}
\end{table}

Using the architectural efficiencies in \autoref{tab:example_effs}, we can calculate the performance portability metrics \pp and \mpp using their definitions \autoref{def:pp} and \autoref{def:mpp}, respectively. 
We use different hardware platform sets, starting with a single hardware platform and gradually adding new platforms to see the change in the performance portability metrics until all six platforms are present. 
In the last two rows, we additionally report the performance portability metric for the platform sets containing all platforms from one specific architectural class. 

The second column contains the proposed heterogeneity score $HS$. 
In the first three rows, the score is $0$ since the platforms $A$, $B$, and $C$ all belong to the same architectural class, completely neglecting the second class. 
After adding platform $D$, the score increases since platforms from both architectural classes are now present. 
The heterogeneity score is $1$ after all platforms are added because then both architectural classes are represented with three platforms each. 
By construction, the last two rows only contain platforms from the same architectural class and, therefore, the heterogeneity score is also $0$. 

For application one, both performance portability metrics have the same characteristics for the first three platforms. 
However, if we add platform D with low architectural efficiency, we see that \pp directly drops to only $\num{0.038}$ since the harmonic mean used in \autoref{def:pp} over proportionally penalizes the newly added low value. 
This does not happen with the \mpp metric, which still has a value of $\num{0.555}$, more closely representing the overall architectural efficiencies. 

For application two, where the architectural efficiencies are of the same magnitude, the performance portability scores of \pp and \mpp are nearly identical. 
Only after adding platform F, which is unsupported by application two, do the scores change: For \pp, the performance portability score immediately drops to zero based on its definition; for \mpp, however, since only supported platforms affect the \mpp score, the value does not change. 
Therefore, looking only at the performance portability values for \mpp, it is impossible to see that application two does not support all hardware platforms. 

For application three, which is not optimized for any of the target platforms, we see small performance portability scores across the board for both \pp and \mpp. 
Again, the values do not diverge much since the architectural efficiencies are also close. 
However, let's look at the final scores when all of the six platforms are added. 
We see another problem with the \pp metric: applications one and three have nearly the same final scores, although they have completely different performance behaviors for the hardware platforms A, B, and C, indicating that information was lost. 

If we follow the advice from \citet{marowka_toward_2021}, we get the performance portability scores shown in the last two rows. 
The scores are nearly identical except for architectural class two and application two, as well as architectural class one and application three. 
However, whether it is sufficient to calculate the scores for the different architectural classes separately depends heavily on the use case. 

All in all, both metrics have their advantages and disadvantages. 
For \pp, it is immediately clear when a hardware platform is unsupported since its score drops to zero; its over-emphasis on small values may be unintuitive and can lose information if only the final scores are reported. 
On the other hand, \mpp tracks the scores in a more intuitive way but also loses information if a hardware platform in the set of all platforms is unsupported for an application. 
In the remainder of this dissertation, I will compare both performance portability metrics \pp and \mpp on real-world problems.
